% ============ 第 3 讲 李嘉图模型（上） ============
\chap{第 3 讲　李嘉图模型（上）}{The Ricardian Model, Part 1}

\textbf{本讲导读}：本讲对应 Lecture 3（The Ricardian Model, Sep.\ 16），共 65 页。内容：贸易原因的分类与各模型的分工；绝对优势 vs.\ 比较优势；李嘉图模型的设定（唯一要素劳动、常数规模报酬）；Home/Foreign 的生产与 PPF；自给自足均衡（劳动市场法）；贸易模式与套利练习；贸易得益（生产法/要素市场法、``迂回生产''）；专业化、就业与真实工资；相对工资与绝对优势。课后手写笔记页（第 64--65 页）给出了 Endowment vs.\ Ricardian 的横向比较，本章``方法地图''以其为骨架。

\secn{3.1 古典模型的路线图与贸易原因的分类}

\textbf{古典贸易模型路线图（第 4--5 页，全课程的分析模板）}：
\begin{itemize}
  \item \term{目标 1：理解贸易的``原因''与``模式''}——根本原因：自给自足下相对价格不同；``迂回''套利：\kw{出口/进口相对价格更低/更高的商品，直到价格收敛}（原文红蓝字）。
  \item \term{路线图第一部分：求自给自足下的相对价格}——
    \begin{itemize}
      \item \kw{PPF 同时也是 CPF}（consumption possibilities frontier，即``预算约束''）：What you produce is what you consume!（蓝字）关键问题：\kw{PPF 的形状与斜率}（荧光高亮）；
      \item 消费者选择：无差异曲线与 PPF 相切（板书草图：$C = Q$）；
      \item 切点斜率 $=$ 自给自足下的相对价格。
    \end{itemize}
  \item \term{目标 2：理解贸易的``后果''}——比较两个情形（贸易 \& 自给自足）；老师手写：贸易均衡中\kw{出口品的相对价格上升}（relative price of exported good $\uparrow$）；荧光高亮：\kw{贸易均衡的相对价格总介于两国自给自足价格之间}；根本变量（Fundamentals，红字）：经济中的相对价格变化。
  \item \term{路线图第二部分：贸易均衡的三步}——
    \begin{itemize}
      \item \kw{生产（Production）}（荧光高亮）：最大化收入 | 要素市场方法；
      \item \term{收入形成（Income formation）}：预算线以某一斜率与 PPF 相切；
      \item \kw{消费（Consumption）}（荧光高亮）：无差异曲线与预算线相切。
    \end{itemize}
\end{itemize}

\textbf{贸易原因的分类（第 7--9 页）}：贸易模型的根本问题——为什么各国会贸易？上一讲：相对价格不等时出现贸易机会——\textbf{为什么不等}？简答：各国在多个维度上不同。国家间的差异构成贸易的原因：
\begin{enumerate}
  \item 各国使用的\kw{技术}不同（荧光高亮）；
  \item 各国\kw{禀赋总量}不同——自然资源、资本、劳动、技能等``\kw{生产要素（factors of production）}''（荧光高亮）；
  \item 各国\kw{偏好（tastes）}不同（荧光高亮）；
  \item 不完全竞争、产品差异化。
\end{enumerate}

\begin{intubox}[各贸易模型的分工（第 9 页，重要导览）]
\textbf{李嘉图（Ricardian）模型}（Lect.\ 3--4）聚焦\kw{技术}差异；\textbf{Heckscher--Ohlin 模型}（Lect.\ 7--8）聚焦\kw{资源/禀赋}差异（荧光高亮）；\textbf{特定要素（Specific Factors）模型}（Lect.\ 6）是技术与禀赋的混合；\textbf{规模经济（Economies of Scale）模型}（Lect.\ 11--13）聚焦产品差异化。（禀赋模型——上一讲——则是最简化的``禀赋差异''基准。）
\end{intubox}

\secn{3.2 绝对优势、比较优势与历史背景}

\begin{defbox}[两种关于技术的论点（第 10 页，两个高亮术语）]
\begin{itemize}
  \item \kw{绝对优势（Absolute advantage）}：一国在生产某商品上拥有\textbf{最好的技术}（best technology）时出口该商品。
  \item \kw{比较优势（Comparative advantage）}：一国\textbf{不需要}比别国生产得更好才能出口；当它生产某商品\textbf{相对自己生产其他商品}做得最好时，就在该商品上有比较优势。
\end{itemize}
\end{defbox}

\begin{intubox}[Samuelson 的评价（第 11 页）]
诺奖得主 Paul Samuelson（1969）曾被数学家 Stanislaw Ulam 挑战``说出社会科学中一个既真又非平凡的命题''。数年后 Samuelson 给出答案：\term{比较优势}。``它逻辑上为真，在数学家面前无需论证；它不平凡，则由成千上万聪明人始终无法自己掌握它、或在被解释后仍无法相信它所证明。''
\end{intubox}

\textbf{李嘉图模型的历史背景（第 12 页）}：
\begin{itemize}
  \item ``\term{重商主义（Mercantilism}''，老师手写中文批注）''：贸易的目标是出口以积累黄金；政策含义：高关税、压低进口、扩大出口。
  \item 老师手写批注：``Gold = Money $\neq$ Wealth''（黄金是货币，不等于财富）；``Net export $\to$ gold $\uparrow \to$ inflation $\to$ $\uparrow$P export good（$\to$ 出口品涨价、竞争力下降）''——休谟的铸币--价格流动机制的课堂口述版。
  \item \term{Smith 与 \kw{Hume}（荧光/彩色高亮）}是早期批评者：出口（卖得比以前贵）与进口（买得比以前便宜）都能获益。
  \item \kw{李嘉图}（粉色高亮）贸易模型：阐明\kw{比较优势}（高亮），以及\textbf{所有国家都从贸易中得益}。
\end{itemize}

\secn{3.3 李嘉图模型设定}

\begin{assumpbox}[李嘉图模型设定（第 14--16 页）]
\begin{enumerate}
  \item \term{两种商品}（老师圈注 $g = W, C$）：小麦（Wheat, W）与布（Cloth, C）。
  \item \term{要素（Factors，老师圈注）}：\kw{只有劳动}（红笔下划线）——没有土地、没有资本；两种商品都只用劳动生产。
  \item \term{技术（Technology，老师圈注）}：常数规模报酬 $Q_g = A_g \cdot L_g$，于是 \kw{$MPL_g = A_g$}（框注）；$MPL_g$（商品 $g$ 的劳动边际产品）是多用一单位劳动的额外产出。等价表达：``\kw{单位劳动投入（unit labor requirement）}''（绿高亮）：生产一单位 $g$ 需要 \kw{$a_g = 1/MPL_g$}（绿高亮）单位劳动。
  \item \term{市场结构}：假设\kw{商品市场与劳动市场都完全竞争}（荧光高亮）——手写批注 ①：\kw{$w = MRPL$}。
  \item \term{对劳动的附加假定}（手写批注 ②）：\kw{劳动在部门间自由流动}（橙笔下划线）；但\term{不能跨国流动}（no migration，橙笔下划线）。
\end{enumerate}
\end{assumpbox}

\textbf{Home 国的数字例子（第 17 页）}：唯一投入是劳动，\kw{要素禀赋}（老师红笔标注``factor endowment''）：圈注 \kw{L = 25 名工人}。Home 一个工人可以生产：4``斗''小麦（$MPL_W = 4$）或 2``尺''布（$MPL_C = 2$）；等价地：生产每单位小麦需 $1/4$ 个工人、每单位布需 $1/2$ 个工人（$a_W = 1/4$，$a_C = 1/2$）。

\textbf{Home 的 PPF（第 18--25 页）}：$L = 25,\; MPL_W = 4,\; MPL_C = 2$。可能的 PPF 形状有三种（图~\ref{fig:ch03-ppfabc}）：
\begin{itemize}
  \item \textbf{(a) 内凹（凸向原点）}：手写注（第 20 页）：``(a) is a weird case that requires \kw{increasing returns to scale by industry}''（行业层面递增规模报酬的怪异情形）——当生产更多 $W$（$Q_W\uparrow, Q_C\downarrow$）时，手写：每增加（边际）一单位 $W$ 需要\kw{更少}资源；每单位资源能生产更多……（手写推导 $\frac{\Delta Q_C}{\Delta Q_W}$ 的变化）。
  \item \textbf{(b) 直线}：手写注（第 21 页）：``(b) is the case under Ricardian model as we compose''——因为：\kw{只有一种要素}、\kw{常数规模报酬}、\kw{完全竞争}。
  \item \textbf{(c) 外凸（凹向原点）}：手写注（第 22 页）：``(c) is the case in future chapters when production requires more than 1 factor''——手写点名：\kw{Specific Factors Model、Heckscher--Ohlin Model}。
\end{itemize}

\begin{figure}[H]
  \centering
  \includegraphics[width=0.5\linewidth]{figs/ch03_fig_01.png}
  \caption{第 19 页：PPF 的三种可能形状——(a) 内凹（行业递增报酬）、(b) 直线（李嘉图：单要素 + CRS + 完全竞争）、(c) 外凸（多要素：特定要素模型与 H-O 模型）。}
  \label{fig:ch03-ppfabc}
\end{figure}

\textbf{PPF 的截距与斜率（第 23--25 页，全部为板书）}：设全部劳动生产布（$L_W = 0, L_C = 25$）：$Q_W = 0, Q_C = 50$，得截距 $A(0, 50)$；全部生产小麦（$L_W = 25, L_C = 0$）：$Q_W = 100, Q_C = 0$，得 $B(100, 0)$（板书作答，见图~\ref{fig:ch03-intercepts}）。

\textbf{斜率的含义}（图~\ref{fig:ch03-intercepts}右）：把 1 单位劳动从布（C）转移到小麦（W）：$\Delta Q_W = 1 \times MPL_W$，$\Delta Q_C = -1 \times MPL_C$，故
\[
\text{Slope} = \frac{\Delta Q_C}{\Delta Q_W} = -\frac{MPL_C}{MPL_W} = -\frac{1}{2}.
\]
即：为多得 1 单位小麦必须放弃的布的数量——\kw{小麦的``机会成本''（opportunity cost）}（橙笔下划线强调）。

\begin{figure}[H]
  \centering
  \begin{minipage}{0.48\linewidth}\centering
    \includegraphics[width=\linewidth]{figs/ch03_fig_02.png}
  \end{minipage}\hfill
  \begin{minipage}{0.48\linewidth}\centering
    \includegraphics[width=\linewidth]{figs/ch03_fig_03.png}
  \end{minipage}
  \caption{左（第 23 页）：PPF 的截距板书——$A(0, 50)$ 与 $B(100, 0)$；右（第 24 页）：斜率板书——把 1 单位劳动从 C 移到 W，Slope $= \frac{\Delta Q_C}{\Delta Q_W} = -\frac{MPL_C}{MPL_W} = -\frac{1}{2}$。}
  \label{fig:ch03-intercepts}
\end{figure}

\secn{3.4 Home 的自给自足均衡}

\textbf{需求面（第 27 页）}：Home 的无差异曲线与 Home PPF（图~\ref{fig:ch03-homeauto}）；无贸易均衡点 $A$ 在 PPF 与无差异曲线 $U_1$ 的切点（$(Q_W, Q_C) = (50, 25)$）。

\begin{figure}[H]
  \centering
  \includegraphics[width=0.55\linewidth]{figs/ch03_fig_04.png}
  \caption{第 27 页：Home 自给自足均衡——无差异曲线与 PPF 相切于 $A$ 点；$U_0 < U_1 < U_2$。}
  \label{fig:ch03-homeauto}
\end{figure}

\textbf{劳动市场均衡（第 28--30 页）}：
\begin{itemize}
  \item 每个行业：工人被雇用到\term{小时工资 $=$ 多生产一小时的价值}为止；后者又等于该小时生产的商品量（MPL）乘以商品价格。手写：\[
  w_g = MRPL_g = \underbrace{MR_g}_{=\,P_g\ (\text{手写红字标注})} \times MPL_g .
  \]
  \item \term{工资应跨行业相等}。手写回答``Why?''：① \kw{劳动完全自由流动}（fully mobile）；② \kw{自给自足均衡下两种商品都被消费}（因此两个行业都必须运转）$\Rightarrow w_W = w_C$（若 $>$ 或 $<$ 则所有劳动流向高工资行业，退化为角点 $L_W = 25, L_C = 0$ 或 $L_W = 0, L_C = 25$，无法两种商品都生产）。
  \item 于是 $P_C \cdot MPL_C = w = P_W \cdot MPL_W$（板书推导）：
  \[
  P_C = \frac{w}{MPL_C}, \qquad P_W = \frac{w}{MPL_W}
  \;\Longrightarrow\;
  \frac{P_W}{P_C} = \frac{MPL_C}{MPL_W}.
  \]
\end{itemize}

\begin{figure}[H]
  \centering
  \includegraphics[width=0.6\linewidth]{figs/ch03_fig_05.png}
  \caption{第 30 页整页板书：劳动市场均衡推导，及手写总结——\kw{相对价格 $=$ ① PPF 的斜率 $=$ ② 小麦的机会成本 $=$ ③ MPL 的反比 $=$ ④ 单位劳动投入之比}。}
  \label{fig:ch03-labormkt}
\end{figure}

\textbf{Foreign 国（第 31--33 页）}：$L^* = 100$ 名工人；一个工人生产 1``斗''小麦（$MPL_W^* = 1$）或 1``尺''布（$MPL_C^* = 1$）。Foreign 的 PPF：截距 $(0, 100)$ 与 $(100, 0)$，\kw{Slope$^* = -\frac{MPL_C^*}{MPL_W^*} = -1$}（板书）。Foreign 自给自足均衡：$\frac{P_W^*}{P_C^*} = 1$（板书作答；图~\ref{fig:ch03-foreignauto}）。红字总结：\kw{相对价格 $=$ MPL 的反比 $=$ PPF 斜率 $=$ 自给自足均衡中预算线的斜率}。

\begin{figure}[H]
  \centering
  \includegraphics[width=0.42\linewidth]{figs/ch03_fig_07.png}
  \caption{第 33 页：Foreign 自给自足均衡——PPF 与无差异曲线切于 $A^*(50, 50)$。}
  \label{fig:ch03-foreignauto}
\end{figure}

\textbf{路线图回顾（第 34 页）}：手写补全——Home 自给自足 $\frac{P_W}{P_C} = \frac{1}{2}$；Foreign 自给自足 $\frac{P_W^*}{P_C^*} = 1$（旁边手写标注 $(\frac{MPL_C}{MPL_W})$、$(\frac{MPL_C^*}{MPL_W^*})$）。路线图 part 1.5：确定贸易模式——比较两国自给自足相对价格（此处即 MPL 的反比，或``单位劳动投入''之比）；\kw{某商品自给自足相对价格更低的国家出口该商品}——这就是``比较优势''的含义（红字）。

\secn{3.5 贸易模式}

\textbf{套利练习（第 36--38 页）}：无贸易时小麦相对价格为 Home $1/2$、Foreign $1$（老师红框圈注）。设想只有你能贸易，从 Home 的 1 单位布开始（板书作答）：
\begin{enumerate}
  \item 在 Home 换成 \kw{2} 单位小麦；
  \item 把 \kw{2} 单位小麦（绿高亮）带到 Foreign；
  \item 在 Foreign 换成 \kw{2} 单位布；
  \item 把 \kw{2} 单位布带回 Home——凭空多出 1 单位布。
\end{enumerate}
板书（第 38 页）：你实际做的事情是\kw{从 Home 出口 W}、\kw{从 Foreign 出口 C}——``\term{按照比较优势贸易}''（Trade according to comparative advantage）！

\textbf{机会成本表与比较优势判定（第 39 页）}：手写笔记（顶部）：``\kw{Relative price} 更技术性；\kw{Opportunity cost} 更哲学、更根本''。

\begin{center}
\begin{tabular}{@{}lcc@{}}
\toprule
 & 1 单位布（C）的机会成本 & 1 单位小麦（W）的机会成本 \\
\midrule
Home & 2 单位小麦 & \kw{1/2 单位布}（高亮） \\
Foreign & \kw{1 单位小麦}（高亮） & 1 单位布 \\
\bottomrule
\end{tabular}
\end{center}

\term{当一国生产某商品的机会成本低于另一国时，该国在该商品上有比较优势（C.A.）}。板书作答：Foreign 在生产 \kw{C（布）}上有比较优势；Home 在生产 \kw{W（小麦）}上有比较优势（红字填空）。

\textbf{走向国际贸易均衡（第 40--41 页）}：随着 Home 持续出口小麦——Home 市场小麦销量 $\downarrow$、Foreign 市场小麦销量 $\uparrow$（板书箭头）：小麦在 Home 的价格 $\uparrow$、在 Foreign 的价格 $\downarrow$；随着 Foreign 持续出口布：布在 Home 的价格 $\downarrow$、在 Foreign 的价格 $\uparrow$；于是\kw{小麦的相对价格在 Home $\uparrow$（从 1/2 起）、在 Foreign $\downarrow$}（板书作答）。

\begin{thmbox}[自由贸易均衡（第 41--42 页）]
``自由贸易均衡''在小麦/布的相对价格在两国相同时实现。求解均衡需要：均衡相对价格 $P_W/P_C$；两国的生产与消费。\kw{自由贸易均衡中小麦的相对价格必介于两国自给自足价格之间}（蓝字强调：between 1/2 and 1）。下一讲求解；本讲先取定 $P_W/P_C \in (1/2, 1)$ 为给定，分析各国的生产与预算线（``消费可能性边界 CPF''）如何变化。
\end{thmbox}

\secn{3.6 贸易得益}

\textbf{路线图回顾（第 44 页）}：part 1（自给自足相对价格：PPF$=$CPF、IC 与 PPF 相切、切点斜率 $=$ 自给自足相对价格）与 part 1.5（比较两国自给自足相对价格）已完成；进入 \kw{part 2：贸易均衡三步}（红字）——\kw{生产}：最大化收入 | 要素市场方法；\term{收入形成}：预算线与 PPF 相切于某一斜率；\kw{消费}：无差异曲线与预算线相切。

\textbf{Home 在 PPF 上哪里生产？（第 45 页，方法 A）}：
\[
\text{Method A:}\quad \max\; I = P_C Q_C + P_W Q_W \quad \text{s.t.}\; \frac{Q_C}{MPL_C} + \frac{Q_W}{MPL_W} = \bar{L} \;(\text{PPF})
\]
手写注释：``\term{iso-income curve（等收入线）}：$I = P_C Q_C + P_W Q_W$ $\Rightarrow$ $P_C Q_C = I - P_W Q_W$ $\Rightarrow$ $Q_C = \frac{I}{P_C} - \frac{P_W}{P_C} Q_W$''（斜率 $-\frac{P_W}{P_C}$，高亮）。图~\ref{fig:ch03-methoda}：等收入线与 PPF 相切的最优点为\kw{角点}——板书圈出横轴端点，标注``\kw{Production（Fully specialize，完全专业化）}''与``CPF / Budget line''。

\begin{figure}[H]
  \centering
  \includegraphics[width=0.62\linewidth]{figs/ch03_fig_08.png}
  \caption{第 45 页：贸易下 Home 的生产决策——等收入线与 PPF 相切于完全专业化角点（$Q_W = 100, Q_C = 0$）；板书标注 CPF（预算线）与``Production (Fully specialize)''。}
  \label{fig:ch03-methoda}
\end{figure}

\textbf{方法 B：观察劳动市场（第 46 页）}：若工人被雇去生产布，挣 $w_C = P_C \cdot MPL_C$；若生产小麦，挣 $w_W = P_W \cdot MPL_W$。\kw{贸易均衡中工人去哪里？}注意贸易均衡满足
\[
\frac{P_W}{P_C} > \text{自给自足相对价格} = \frac{MPL_C}{MPL_W}
\;\Longrightarrow\;
\text{手写：}\; \underbrace{P_W \cdot MPL_W}_{\tilde{w}_W} > \underbrace{P_C \cdot MPL_C}_{\tilde{w}_C}
\]
——小麦部门工资更高，全部劳动流向小麦部门：Home 完全专业化于小麦。

\textbf{两国的贸易得益（第 47--48 页）}：Home 能达到\kw{更高的效用水平}（高亮）——\term{贸易得益为正（Positive gains from trade）}（粉框强调）！手写批注：判断依据可以是``指定效用函数''或``\kw{显示偏好（Revealed Preference）}''。Foreign 同样能达到更高效用水平——\kw{两国贸易得益均为正}（Positive gains from trade for both countries）。图~\ref{fig:ch03-foreigngain}：Foreign 在 $B^*(0,100)$ 完全专业化生产布，按世界价格线（斜率 $-\frac{2}{3}$）出口 40 码布、进口 60 蒲式耳小麦，消费点 $C^*$ 达到 $U_2^* > U_1^*$。手写总结四点：① $\frac{P_W}{P_C} < (\frac{P_W^*}{P_C^*})^{autarky}$；② 最优消费：专业化于 C；③ 更高效用；④ \kw{选择集扩张（choice set expansion）}。

\begin{figure}[H]
  \centering
  \includegraphics[width=0.5\linewidth]{figs/ch03_fig_09.png}
  \caption{第 48 页：Foreign 的贸易得益——生产点 $B^*$、消费点 $C^*$、世界价格线斜率 $-2/3$（介于 $-1$ 与 $-1/2$ 之间）；出口 40 码布、进口 60 蒲式耳小麦。}
  \label{fig:ch03-foreigngain}
\end{figure}

\textbf{理解贸易得益：``迂回生产''（第 49--51 页）}：
\begin{itemize}
  \item 贸易好比\term{布的 MPL 提升}。如何理解？Home 通过贸易\term{以``迂回方式''（roundabout way）生产布}：
  \begin{itemize}
    \item 直接方式：1 单位劳动 $\to$ \kw{2}（$MPL_C$）单位布（板书作答）；
    \item 迂回方式：1 单位劳动 $\to$ \kw{4}（$MPL_W$）单位小麦（板书）$\to$ 按贸易均衡相对价格换布。板书设均衡相对价格为 \kw{2/3}（圈注；小麦的相对价格，介于 1/2 与 1 之间）：4 单位小麦 $\to$ $4 \times \frac{2}{3} = \frac{8}{3}$ 单位布。
  \end{itemize}
  \item 由于 $\frac{8}{3} > 2$，\redkw{迂回方式比直接方式更有效率（more efficient）}；
  \item \kw{贸易就像布的 MPL 从 2 提高到 $\frac{8}{3}$}（板书填空）；
  \item 用 $MPL_C = \frac{8}{3}$ 重算 PPF $\Rightarrow$ \kw{预算集（CPF）扩张}。
\end{itemize}

\textbf{贸易得益的三重类比（第 58 页，对照 GDP 的三种定义）}：
\begin{center}
\begin{tabular}{@{}lll@{}}
\toprule
GDP 的定义 & 贸易得益的对应表述 \\
\midrule
支出之和（手写 $C + I + G + NX$） & CPF 扩张——代表性主体能\kw{消费}更多（``overall''，手写） \\
生产之和（手写 $\sum_{i \in \text{final}} P_i Y_i$） & ``迂回''生产——进口品的等价\kw{生产率}更高 \\
收入之和（手写 $wL + rK + \cdots$） & \kw{真实工资}上升（手写注：potentially grosso modo，粗略地） \\
\bottomrule
\end{tabular}
\end{center}

\secn{3.7 专业化、就业与工资}

\textbf{完全专业化（第 53--54 页）}：贸易均衡中，各国倾向于\term{专业化生产其有比较优势的产品}。\kw{Home 生产小麦}（粉框）：$w = P_W \cdot MPL_W$（板书打勾），与 $P_C \cdot MPL_C$ 比较（手写：$>$）；\kw{Foreign 生产布}（粉框）：$w^* = P_C \cdot MPL_C^*$（手写：$> P_W \cdot MPL_W^*$）。

机制（第 54 页）：Home 专业化于小麦，\kw{布产业在 Home 消失}——只要 $\frac{P_W}{P_C} > \frac{MPL_C}{MPL_W}$（贸易相对价格高于自给自足），就有 $P_W \cdot MPL_W > P_C \cdot MPL_C$：布产业无法吸引任何工人（工资低于小麦业）。

\textbf{真实工资（第 55--57 页）}：\term{以商品 $g$ 计的真实工资} $= \frac{w}{P_g}$（红字定义）。工人的福利用两种商品的真实工资衡量。

\begin{figure}[H]
  \centering
  \includegraphics[width=0.8\linewidth]{figs/ch03_fig_10.png}
  \caption{第 56 页板书：Home 的真实工资——贸易后 $w/P_W = MPL_W$ 不变（手写``不变''），$w/P_C = \frac{P_W}{P_C} \cdot MPL_W$ 提高（手写``提高''）；手写注：$\frac{P_W}{P_C} > (\frac{P_W}{P_C})^{autarky} = \frac{MPL_C}{MPL_W}$ $\Rightarrow$ \kw{Workers' welfare? Improve!}（工人福利改善）。}
  \label{fig:ch03-realwageH}
\end{figure}

\begin{figure}[H]
  \centering
  \includegraphics[width=0.8\linewidth]{figs/ch03_fig_11.png}
  \caption{第 57 页板书（课后练习）：Foreign 的真实工资——贸易后 $w^*/P_C = MPL_C^*$ 不变（手写 no change），$w^*/P_W = \frac{P_C}{P_W} \cdot MPL_C^*$ 增加（手写 increase）$\Rightarrow$ \kw{Worker welfare in F improve}（手写）。}
  \label{fig:ch03-realwageF}
\end{figure}

\begin{thmbox}[李嘉图模型的分配结论]
贸易后：出口部门\textbf{以本部门商品计的真实工资不变}（$w/P_W = MPL_W$），以进口品计的真实工资上升（$w/P_C = \frac{P_W}{P_C} MPL_W \uparrow$，因为 $P_W/P_C$ 上升）——\kw{工人的福利（唯一要素的所有者）确定改善}。本模型中不存在输家：唯一要素劳动可跨部门流动，人人都在扩张的出口部门就业。
\end{thmbox}

\textbf{相对工资（Relative Wages，第 59--61 页）}：定义 $\frac{w}{w^*}$（板书）。Home 生产小麦 $w = P_W \cdot MPL_W$；Foreign 生产布 $w^* = P_C \cdot MPL_C^*$。比较两国工资有两层含义：
\begin{enumerate}
  \item \term{工人的相对福利（Relative welfare of workers）}。相对工资的范围：自给自足价格满足 $\frac{MPL_C}{MPL_W} \le \frac{P_W}{P_C} \le \frac{MPL_C^*}{MPL_W^*}$（板书：$\frac{1}{2} \le \frac{P_W}{P_C} \le 1$）。于是
  \[
  \frac{w}{w^*} = \frac{P_W}{P_C} \times \frac{MPL_W}{MPL_C^*} \in \left[ \frac{MPL_C}{MPL_C^*},\; \frac{MPL_W}{MPL_W^*} \right]
  \]
  （板书代入本例数字：$[2/1, 4/1]$，即 $w/w^* \in [2, 4]$；手写注：区间端点对应``\kw{Absolute Advantage}（绝对优势）''——板书框出 $\frac{MPL_C}{MPL_C^*}$ 与 $\frac{MPL_W}{MPL_W^*}$）。相对工资越高，Home 工人相对福利越高。本例中 Home 在\textbf{两种}商品上都有绝对优势（$MPL_W = 4 > 1$，$MPL_C = 2 > 1$），故 Home 的相对工资恒大于 1。
  \item \term{劳动的相对成本（Relative cost of labor）}：从相对工资可以``反向工程''出专业化模式（第 60 页表）：
\end{enumerate}

\begin{center}
\begin{tabular}{@{}lcc@{}}
\toprule
生产一单位的劳动成本 & 小麦 & 布 \\
\midrule
Home & $\frac{w}{MPL_W} = \frac{1}{4} w$（红字） & $\frac{w}{MPL_C} = \frac{1}{2} w$（红字） \\
Foreign & $\frac{w^*}{MPL_W^*} = w^*$ & $\frac{w^*}{MPL_C^*} = w^*$ \\
Home 生产该商品当 & $\frac{1}{4} w \le w^*$（即 $\frac{w}{w^*} \le 4$） & $\frac{1}{2} w \le w^*$（即 $\frac{w}{w^*} \le 2$） \\
Foreign 生产该商品当 & $w^* \le \frac{1}{4} w$（即 $\frac{w}{w^*} \ge 4$） & $w^* \le \frac{1}{2} w$（即 $\frac{w}{w^*} \ge 2$） \\
\bottomrule
\end{tabular}
\end{center}

\begin{figure}[H]
  \centering
  \includegraphics[width=0.85\linewidth]{figs/ch03_fig_12.png}
  \caption{第 61 页板书：以相对工资 $w/w^*$ 为轴``反向工程''专业化模式——$w/w^* < 2$：F 无产业可竞争（不可均衡，板书打叉）；$w/w^* = 2$ 与 $= 4$：不完全专业化（板书``???''）；$2 < w/w^* < 4$：\kw{完全专业化}（H 生产 W、F 生产 C）；$w/w^* > 4$：H 无产业（不可均衡，打叉）。}
  \label{fig:ch03-relwage}
\end{figure}

\begin{intubox}[讨论：社会计划者视角 vs.\ 微观主体视角（第 62 页）]
模型中我们用\term{比较优势}讨论``贸易原因''——基于整个经济的 PPF 与机会成本（老师手写批注``\kw{Social Planner's View}''）。而现实中与任何做贸易的企业交谈——他们考虑的是自己企业的投入产出选择（手写批注``\kw{Micro-agents' View}''），``有利可图就贸易''。\redkw{我们的理论错了吗？若没错，两种视角如何调和？}——老师要求同学准备自己的答案，上半学期总结课回顾。
\end{intubox}

\textbf{下一讲预告（第 63 页）}：求解贸易均衡；李嘉图模型的比较静态；李嘉图模型的扩展（劳动不可流动；非贸易品）。

\secn{3.8 模型构建方法地图（据课后手写笔记整理）}

\begin{mapbox}[Endowment 与 Ricardo 模型的横向比较（课后手写笔记第 64--65 页）]
\textbf{1. 共同的 insights}：
\begin{itemize}
  \item \term{贸易原因}：自给自足下相对价格差异 $\Rightarrow$ 套利；
  \item \term{贸易模式}：出口自给自足下相对价格低的产品；
  \item \term{贸易平衡}：同一产品在两国相对价格趋同；
  \item 核心思维实验：``\kw{roundabout trade}（迂回贸易）''。
\end{itemize}
\textbf{2. Autarky 均衡}（手写对照表，见图~\ref{fig:ch03-note1}）：
\begin{center}
\begin{tabular}{@{}p{0.22\linewidth}p{0.34\linewidth}p{0.34\linewidth}@{}}
\toprule
 & Endowment & Ricardo \\
\midrule
PPF & \term{点}（禀赋） & \term{直线}，slope $= -\frac{MPL_C}{MPL_W} = -\frac{P_W}{P_C}$ \\
Autarky 均衡条件 & 生产 $=$ 消费：$(Q_W, Q_C) = (C_W, C_C)$（``自给自足''） & 同左（生产点与消费点重合） \\
相对价格 & 由禀赋相对丰裕程度（P指数）决定 & 由两商品 $MPL_g$ 之 inverse ratio 得出 \\
要素市场 & 不存在 & 劳动力自由流动：$w = P_C \cdot MPL_C = P_W \cdot MPL_W$ \\
\bottomrule
\end{tabular}
\end{center}
\textbf{3. 贸易对经济之影响}（手写注：对 Home 而言，技术上视为 $\frac{P_W}{P_C}$ 增加，即\kw{出口品的相对价格提升}；见图~\ref{fig:ch03-note2}）：
\begin{center}
\begin{tabular}{@{}p{0.22\linewidth}p{0.34\linewidth}p{0.34\linewidth}@{}}
\toprule
 & Endowment & Ricardian \\
\midrule
生产决策 & 无 & 在具有比较优势的部门专业化：$Q^W = L \cdot MPL_W$，$Q^C = 0$ \\
收入 & \multicolumn{2}{l}{卖掉所有商品：$I = P_W Q_W + P_C Q_C$} \\
CPF 图示 / 消费选择 & 消费点沿 CPF 移动（出口品减少、进口品增加） & 同左（CPF 从禀赋点/专业化点出发） \\
Gains from trade & ① 效用明显提升；② 消费选择扩展（``支出法''） & ① 效用明显提升；② 消费选择扩展（``支出法''）；③ \kw{生产效率提升} $MPL_C \uparrow$（``Roundabout''间接生产）；④ \kw{劳动真实收入提升}：$w/P_C \uparrow$、$w/P_W$ 不变（``收入法''） \\
\bottomrule
\end{tabular}
\end{center}
（手写注：Endowment 中没有生产或要素收入——所以只有 ①② 两条渠道；Ricardo 多出 ③④ 两条。）
\end{mapbox}

\begin{figure}[H]
  \centering
  \includegraphics[width=0.6\linewidth]{figs/ch03_fig_13.png}
  \caption{课后手写笔记（一）：Endowment 与 Ricardo 横向比较——共同 insights、自给自足均衡对照表（PPF、均衡条件、相对价格、要素市场）。}
  \label{fig:ch03-note1}
\end{figure}

\begin{figure}[H]
  \centering
  \includegraphics[width=0.55\linewidth]{figs/ch03_fig_14.png}
  \caption{课后手写笔记（二）：贸易对经济之影响——CPF 图示与贸易得益的四条渠道（① 效用提升 ② 消费选择扩展 ③ 生产效率提升 ④ 劳动真实收入提升）。}
  \label{fig:ch03-note2}
\end{figure}

\begin{sumbox}[本讲总结：通用结论与方法]
\begin{enumerate}
  \item \term{贸易原因分类}：技术差异（李嘉图）、禀赋差异（H-O / 禀赋模型）、偏好差异、规模经济与产品差异化——每个模型只强调一个维度。
  \item \term{比较优势 $\neq$ 绝对优势}：出口模式由\textbf{相对}机会成本决定；即使一国在所有商品上都有绝对优势（本例 Home：$MPL_W = 4 > 1$，$MPL_C = 2 > 1$），两国仍都从贸易中获益——李嘉图模型的``非平凡''之处（Samuelson）。
  \item \term{李嘉图设定}：单要素（劳动）、CRS（$MPL_g = A_g$ 常数）、完全竞争、劳动国内自由流动 / 国际不可流动；PPF 为直线，斜率 $= -\frac{MPL_C}{MPL_W} = -\frac{1}{2}$（本例）。
  \item \term{自给自足相对价格四重身份}（板书）：$P_W/P_C = $ ① PPF 斜率 $=$ ② 小麦的机会成本 $=$ ③ MPL 反比（$\frac{MPL_C}{MPL_W}$）$=$ ④ 单位劳动投入之比（$\frac{a_W}{a_C}$）。劳动市场均衡条件 $w = P_g \cdot MPL_g$ 且跨行业相等。
  \item \term{贸易模式}：出口自给自足相对价格更低的商品（Home 出口小麦、Foreign 出口布）；世界相对价格介于两国自给自足价格之间（$(1/2, 1)$）。
  \item \term{贸易得益的四种表述}：① 效用提升（显示偏好/指定效用函数）；② CPF 扩张（支出法）；③ ``迂回生产''：贸易像把进口品部门的 MPL 从 2 提高到 $MPL_W \times \frac{P_W}{P_C} = \frac{8}{3}$（生产法）；④ 出口部门真实工资不变、进口品计的真实工资上升（收入法）。
  \item \term{完全专业化}：只要贸易相对价格严格高于自给自足相对价格，出口部门工资更高，劳动全部流向出口部门，进口部门消失；$w/w^* \in [\frac{MPL_C}{MPL_C^*}, \frac{MPL_W}{MPL_W^*}]$ 之外的模式不可均衡，边界上出现不完全专业化。
  \item \term{相对工资的两重含义}：工人相对福利；劳动相对成本（可反向工程专业化模式）。李嘉图模型中无输家：唯一要素劳动完全流动。
\end{enumerate}
\end{sumbox}
